% GENERATED FILE. Edit canonical lessons, then run npm run generate:books.
% canonical-source-hash: fnv1a64:c1f4a8bd4d2dd0e3
% canonical-lessons: ZH-C87-meiyou, ZH-W87-mei, ZH-C87-bu-mei, ZH-C87-chuang, ZH-C87-deng, ZH-C87-yaoshi, ZH-C87-bingxiang, ZH-C87-practice

\chapter{Is There Any?}
\label{ch:zh-is-there-any}
\hlchaptermodality{87}

\begin{chapteropening}
\textbf{By the end of this chapter:} \emph{I can say what there is and is not, and that something did not happen.}

Saying what there is and is not with \zh{没有}, telling \zh{不} from \zh{没}, and naming a bed, a light, a key and a fridge, built on one new written character: \zh{没}.
\end{chapteropening}

\section[\texorpdfstring{\zh{没有} (méiyǒu)}{没有 (méiyǒu)}]{\zh{没有} (méiyǒu) — not have; there is no}

\label{lesson:ZH-C87-meiyou}



\begin{quote}
A new chapter, and it opens with the other half of a word you have used since the first chapters.

\begin{itemize}
\raggedright
  \item \emph{Recall:} say \zh{你在哪儿} — \emph{nǐ zài nǎr}, and answer it — \textbf{R1}, the chapter you have finished
  \item \emph{Recall:} say \emph{xǐshǒujiān}, the toilet — \textbf{R1}
  \item \emph{Recall:} say \emph{wǒ yǒu shíjiān}, I have time
\end{itemize}
\end{quote}

\subsection*{You'll want to know: \zh{没有}}
\textbf{\zh{没有}} — \emph{méiyǒu} — \textquotedblleft{}not have; there is no\textquotedblright{}.

\begin{quote}
\textbf{\zh{我没有手机}} \emph{wǒ méiyǒu shǒujī} — I don't have a phone.
\end{quote}

\begin{quote}
\textbf{\zh{这儿没有人}} \emph{zhèr méiyǒu rén} — There's nobody here.
\end{quote}

The next lesson puts the new character in your hand.

\begin{cousinweb}
\textbf{\zh{没有}} (\emph{méiyǒu}) is built from \zh{没} 'not' and \zh{有} 'to have'. \zh{有} is inherited from Middle Chinese \emph{hjuwX} (Guangyun, in Baxter's transcription). \zh{没} is there as \emph{mwot}, glossed 'to sink'; the Guangyun does not give 'not', and the regular sound change from \emph{mwot} gives mò, not méi, so how \zh{没} came to mean 'not' and to be said méi is not settled.
\end{cousinweb}

\begin{grammarlens}[title={What you've built}]
\zh{有} has its own negative, and it is \zh{没}, not \zh{不}: \zh{没有}. Both jobs of \zh{有} go with it, 'to have' and 'there is': \zh{我没有手机}, \zh{这儿没有人}.
\end{grammarlens}

\subsection*{Your turn}
\begin{itemize}
\raggedright
  \item \emph{Say it:} \emph{méiyǒu}
  \item \emph{Say it:} \emph{méiyǒu}, once more
  \item \emph{Say it:} \zh{我没有手机} — \emph{wǒ méiyǒu shǒujī} — I don't have a phone.
  \item \emph{Say it:} \zh{这儿没有人} — \emph{zhèr méiyǒu rén} — There's nobody here.
  \item \emph{Recall:} say \emph{wǒ yǒu shíjiān}, then turn it round: \emph{wǒ méiyǒu shíjiān}
\end{itemize}

\begin{tcolorbox}[breakable,colback=teal!4,colframe=teal!35!black,title={Before you move on}]
What does \zh{没有} mean? (\textquotedblleft{}Not have; there is no\textquotedblright{}.) Which negative goes with \zh{有}? (\textbf{\zh{没}}.)
\end{tcolorbox}
\section[\texorpdfstring{\zh{没} (méi)}{没 (méi)}]{\zh{没} — water, a small roof, and \zh{又}}

\label{lesson:ZH-W87-mei}



\begin{quote}
The word first, then the character inside it.

\begin{itemize}
\raggedright
  \item \emph{Recall:} say \emph{méiyǒu} — \textbf{R1}, one lesson back, before it decays
\end{itemize}

\begin{quote}
\zh{汉} \zh{对}
\end{quote}

\emph{Read it:} each one aloud, with nothing said first — \emph{hàn}, \emph{duì}
\end{quote}

\subsection*{Script — water, a small roof, and \zh{又}}
\hlblockfigure{\detokenize{figures/ZH-W87-mei-filmstrip.pdf}}{How \zh{没} is written, stroke by stroke}

\begin{quote}
\zh{没}
\end{quote}

\textbf{Three pieces, left before right, top before bottom.}

\textbf{The water piece \zh{氵}, three strokes}, as \zh{汉} opens: two dots, then a stroke rising to the right.

\textbf{A small roof, two strokes}: a short left-falling stroke, then a horizontal that turns down and runs out to the right.

\textbf{\zh{又}, two strokes}, under the roof, as \zh{对} opens: across and falling to the left, then the long stroke falling to the right.

3 + 2 + 2 = seven.

\subsection*{Your turn}
\begin{itemize}
\raggedright
  \item \emph{Write it:} \zh{没}, all seven, with the model in front of you
  \item \emph{Cover:} write it again from memory, and count the pieces, not the strokes
\end{itemize}

\begin{tcolorbox}[breakable,colback=teal!4,colframe=teal!35!black,title={Before you move on}]
How many strokes in \zh{没}? (\textbf{Seven}.) Which character opens with the same water piece? (\textbf{\zh{汉}}.) Which shape sits under the roof? (\textbf{\zh{又}, as in \zh{对}}.)

Source: \href{https://raw.githubusercontent.com/chanind/hanzi-writer-data/68d10a4b21150cae5e1ebbd223eed289cf32d90c/data/\%E6\%B2\%A1.json}{Hanzi Writer Data \zh{没}, pinned stroke paths and medians}.
\end{tcolorbox}
\section[\texorpdfstring{\zh{不} — \zh{没} (bù — méi)}{不 — 没 (bù — méi)}]{\zh{不} — \zh{没}: two words for not}

\label{lesson:ZH-C87-bu-mei}



\begin{quote}
Two words for 'not', and this lesson sorts them.

\begin{itemize}
\raggedright
  \item \emph{Read it:} \zh{没有}, aloud — \emph{méiyǒu}
  \item \emph{Recall:} say \zh{不} — \emph{bù}, not
  \item \emph{Recall:} say \zh{好不好} — \emph{hǎo bu hǎo}, is that all right?
\end{itemize}
\end{quote}

\begin{grammarlens}[title={Grammar Lens: \zh{不} for now and for habit, \zh{没} for what did not happen}]
\begin{quote}
\textbf{\zh{不} says no about now, about habit, and about what you will or won't do. \zh{没} says something did not happen.}
\end{quote}

\begin{quote}
\textbf{\zh{我不去}} \emph{wǒ bú qù} — I'm not going; I don't go. \textbf{\zh{我没去}} \emph{wǒ méi qù} — I didn't go.
\end{quote}

Same verb, same person; the negative alone moves the sentence from a decision to a report. \zh{不} before the fall of \zh{去} is \emph{bú}, as in \emph{bú shì}.

And \zh{有} takes \zh{没} every time, which is how the question works:

\begin{quote}
\textbf{\zh{你有没有手机}} \emph{nǐ yǒu méiyǒu shǒujī} — Do you have a phone?
\end{quote}

It is the same yes-or-no shape as \zh{好不好}: the verb, its negative, the verb again. The answer is \zh{有} or \zh{没有}.
\end{grammarlens}

\subsection*{Your turn}
Say these aloud:

\begin{itemize}
\raggedright
  \item \zh{我不去} — \emph{wǒ bú qù} — I'm not going.
  \item \zh{我没去} — \emph{wǒ méi qù} — I didn't go.
  \item \zh{你有没有手机} — \emph{nǐ yǒu méiyǒu shǒujī} — Do you have a phone?
  \item \zh{没有} — \emph{méiyǒu} — No, I don't.
\end{itemize}

\begin{tcolorbox}[breakable,colback=teal!4,colframe=teal!35!black,title={Before you move on}]
Which negative says something did not happen? (\textbf{\zh{没}}.) Which one says you won't? (\textbf{\zh{不}}.) How do you ask whether somebody has a phone? (\textbf{\zh{你有没有手机}}.)
\end{tcolorbox}
\section[chuáng]{chuáng — bed}

\label{lesson:ZH-C87-chuang}



\begin{quote}
Before the new one: the two negatives, then two words for home.

\begin{itemize}
\raggedright
  \item \emph{Recall:} say \zh{我没去} — \emph{wǒ méi qù}, I didn't go
  \item \emph{Recall:} say \emph{fángjiān}, a room — \textbf{R4}, recognition at distance
  \item \emph{Recall:} say \emph{qǐchuáng}, to get up — \textbf{R3}, durable
\end{itemize}
\end{quote}

\subsection*{You'll want to know: chuáng}
\emph{chuáng} — \textquotedblleft{}bed\textquotedblright{}.

\begin{quote}
\emph{fángjiān lǐ yǒu chuáng ma} — Is there a bed in the room?
\end{quote}

Its character comes later; for now it is a word for the ear.

\begin{cousinweb}
\emph{chuáng} is inherited from Middle Chinese \emph{dzrjang} (Guangyun, in Baxter's transcription), glossed 'a bed of slats'. It is the chuáng of \emph{qǐchuáng}: to rise from the bed.
\end{cousinweb}

\begin{grammarlens}[title={What you've built}]
\emph{fángjiān lǐ yǒu chuáng ma} is \zh{有} in its 'there is' job: the place first, then \zh{有}, then the thing. The answer is \zh{有} or \zh{没有}.
\end{grammarlens}

\subsection*{Your turn}
\begin{itemize}
\raggedright
  \item \emph{Say it:} \emph{chuáng}
  \item \emph{Say it:} \emph{chuáng}, once more
  \item \emph{Say it:} \emph{fángjiān lǐ yǒu chuáng ma} — Is there a bed in the room?
  \item \emph{Say it:} \zh{没有} — \emph{méiyǒu} — No, there isn't.
  \item \emph{Recall:} say \emph{qǐchuáng}, then \emph{chuáng} again
\end{itemize}

\begin{tcolorbox}[breakable,colback=teal!4,colframe=teal!35!black,title={Before you move on}]
What does \emph{chuáng} mean? (\textquotedblleft{}Bed\textquotedblright{}.) Which word you know ends with it? (\emph{\textbf{qǐchuáng}}, to get up.)
\end{tcolorbox}
\section[dēng]{dēng — lamp; light}

\label{lesson:ZH-C87-deng}



\begin{quote}
Before the new one: last lesson's word, then the two verbs that go with this one.

\begin{itemize}
\raggedright
  \item \emph{Recall:} say \emph{chuáng}, bed — \textbf{R1}, one lesson back
  \item \emph{Recall:} say \emph{kāi}, to switch on — \textbf{R3}, durable
  \item \emph{Recall:} say \emph{guān}, to switch off — \textbf{R3}
\end{itemize}
\end{quote}

\subsection*{You'll want to know: dēng}
\emph{dēng} — \textquotedblleft{}lamp; light\textquotedblright{}.

\begin{quote}
\emph{kāi dēng} — Switch on the light.
\end{quote}

\begin{quote}
\emph{guān dēng} — Switch off the light.
\end{quote}

Its character comes later; for now it is a word for the ear.

\begin{cousinweb}
\emph{dēng} is inherited from Middle Chinese \emph{tong} (Guangyun, in Baxter's transcription), glossed 'a lamp'.
\end{cousinweb}

\begin{grammarlens}[title={What you've built}]
\emph{kāi} and \emph{guān} open and shut doors, and the same two verbs switch a light on and off.
\end{grammarlens}

\subsection*{Your turn}
\begin{itemize}
\raggedright
  \item \emph{Say it:} \emph{dēng}
  \item \emph{Say it:} \emph{dēng}, once more
  \item \emph{Say it:} \emph{kāi dēng} — Switch on the light.
  \item \emph{Say it:} \emph{guān dēng} — Switch off the light.
  \item \emph{Recall:} say \emph{chuáng}, then \emph{dēng} again
\end{itemize}

\begin{tcolorbox}[breakable,colback=teal!4,colframe=teal!35!black,title={Before you move on}]
What does \emph{dēng} mean? (\textquotedblleft{}Lamp; light\textquotedblright{}.) Switch it on. (\emph{kāi dēng}.)
\end{tcolorbox}
\section[yàoshi]{yàoshi — key}

\label{lesson:ZH-C87-yaoshi}



\begin{quote}
Before the new one: last lesson's word, then two words for coming in.

\begin{itemize}
\raggedright
  \item \emph{Recall:} say \emph{dēng}, a light — \textbf{R1}, one lesson back
  \item \emph{Recall:} say \zh{门} — \emph{mén}, a door
  \item \emph{Recall:} say \zh{进} — \emph{jìn}, to come in
\end{itemize}
\end{quote}

\subsection*{You'll want to know: yàoshi}
\emph{yàoshi} — \textquotedblleft{}key\textquotedblright{}.

\begin{quote}
\emph{wǒ méiyǒu yàoshi} — I don't have a key.
\end{quote}

Its characters come later; for now it is a word for the ear.

\begin{cousinweb}
\emph{yàoshi} is built from two words inherited from Middle Chinese, \emph{yak} and \emph{dzye} (Guangyun, in Baxter's transcription), glossed 'a bolt' and 'a spoon'. The Guangyun does not give the pair as 'key', so how the two came together is not settled.
\end{cousinweb}

\begin{grammarlens}[title={What you've built}]
\emph{wǒ méiyǒu yàoshi}: \zh{没有} in its 'not have' job, with nothing between it and the thing.
\end{grammarlens}

\subsection*{Your turn}
\begin{itemize}
\raggedright
  \item \emph{Say it:} \emph{yàoshi}
  \item \emph{Say it:} \emph{yàoshi}, once more
  \item \emph{Say it:} \emph{wǒ méiyǒu yàoshi} — I don't have a key.
  \item \emph{Say it:} \emph{mén}, then \emph{yàoshi}
  \item \emph{Recall:} say \emph{dēng}, then \emph{yàoshi} again
\end{itemize}

\begin{tcolorbox}[breakable,colback=teal!4,colframe=teal!35!black,title={Before you move on}]
What does \emph{yàoshi} mean? (\textquotedblleft{}Key\textquotedblright{}.) Say you don't have one. (\emph{wǒ méiyǒu yàoshi}.)
\end{tcolorbox}
\section[bīngxiāng]{bīngxiāng — fridge}

\label{lesson:ZH-C87-bingxiang}



\begin{quote}
Before the new one: last lesson's word, then a drink, then a word for hurrying.

\begin{itemize}
\raggedright
  \item \emph{Recall:} say \emph{yàoshi}, a key — \textbf{R1}, one lesson back
  \item \emph{Recall:} say \emph{niúnǎi}, milk — \textbf{R4}, recognition at distance
  \item \emph{Recall:} say \zh{马上} — \emph{mǎshàng}, right away — \textbf{R2}
\end{itemize}
\end{quote}

\subsection*{You'll want to know: bīngxiāng}
\emph{bīngxiāng} — \textquotedblleft{}fridge; refrigerator\textquotedblright{}.

\begin{quote}
\emph{bīngxiāng lǐ yǒu niúnǎi} — There's milk in the fridge.
\end{quote}

Its characters come later; for now it is a word for the ear.

\begin{cousinweb}
\emph{bīngxiāng} is built from bīng 'ice' and xiāng 'a box', both inherited from Middle Chinese, \emph{ping} and \emph{sjang} (Guangyun, in Baxter's transcription), glossed 'frozen water' and 'a box; a chest'. The pair names a machine the Guangyun never saw: the compound is modern.
\end{cousinweb}

\begin{grammarlens}[title={What you've built}]
The place first, then \zh{有}, then the thing: \emph{bīngxiāng lǐ yǒu niúnǎi}. The \zh{里} after \emph{bīngxiāng} is the one you write, and it means inside.
\end{grammarlens}

\subsection*{Your turn}
\begin{itemize}
\raggedright
  \item \emph{Say it:} \emph{bīngxiāng}
  \item \emph{Say it:} \emph{bīngxiāng}, once more
  \item \emph{Say it:} \emph{bīngxiāng lǐ yǒu niúnǎi} — There's milk in the fridge.
  \item \emph{Say it:} \emph{bīngxiāng lǐ méiyǒu shuǐ} — There's no water in the fridge.
  \item \emph{Recall:} say \emph{yàoshi} and \emph{mǎshàng}, then \emph{bīngxiāng} again
\end{itemize}

\begin{tcolorbox}[breakable,colback=teal!4,colframe=teal!35!black,title={Before you move on}]
What does \emph{bīngxiāng} mean? (\textquotedblleft{}Fridge\textquotedblright{}.) Say there is milk in it. (\emph{bīngxiāng lǐ yǒu niúnǎi}.)
\end{tcolorbox}
\section[Practice]{\zh{你有没有} — \zh{没有}}

\label{lesson:ZH-C87-practice}



\begin{quote}
The chapter closes here, with a guest shown to a room.

\begin{itemize}
\raggedright
  \item \emph{Recall:} say \emph{bīngxiāng}, a fridge — \textbf{R1}, one lesson back
  \item \emph{Recall:} say \emph{chídào}, to be late — \textbf{R2}
  \item \emph{Recall:} say \emph{láibují}, to run out of time — \textbf{R2}
\end{itemize}
\end{quote}

\subsection*{The exchange}
\begin{quote}
\textbf{A:} \emph{zhè shì nǐ de fángjiān} — This is your room.
\end{quote}

\begin{quote}
\textbf{B:} \zh{谢谢}
\end{quote}

\begin{quote}
\emph{xièxie} — Thanks.
\end{quote}

\begin{quote}
\textbf{B:} \emph{fángjiān lǐ yǒu dēng ma} — Is there a light in the room?
\end{quote}

\begin{quote}
\textbf{A:} \emph{yǒu, zài mén pángbiān} — Yes, by the door.
\end{quote}

\begin{quote}
\textbf{B:} \emph{nǐ yǒu méiyǒu yàoshi} — Have you got a key?
\end{quote}

\begin{quote}
\textbf{A:} \zh{没有}
\end{quote}

\begin{quote}
\emph{méiyǒu} — No, I haven't.
\end{quote}

\begin{quote}
\textbf{A:} \emph{bīngxiāng lǐ yǒu niúnǎi} — There's milk in the fridge.
\end{quote}

\begin{quote}
\textbf{B:} \zh{我不喝}
\end{quote}

\begin{quote}
\emph{wǒ bù hē} — I don't drink it.
\end{quote}

Two kinds of no in one room: \zh{没有}, there isn't any, and \zh{我不喝}, I don't drink it. The first reports, the second refuses.

\begin{tcolorbox}[breakable,skin=enhanced,colback=orange!6,colframe=orange!45!black,boxrule=0.5pt,arc=1mm,left=6pt,right=6pt,top=4pt,bottom=4pt,fonttitle=\bfseries,title={Read this}]
\emph{Read it:} the exchange through once without stopping, then answer the question below
\end{tcolorbox}

\begin{grammarlens}[title={What you've built — four things, and what is missing}]
A bed, a light, a key and a fridge, and two ways to say a thing is not there:

\begin{itemize}
\raggedright
  \item \emph{nǐ yǒu méiyǒu yàoshi} — \textbf{\zh{没有}} \emph{méiyǒu}
  \item \emph{fángjiān lǐ yǒu chuáng ma} — \textbf{\zh{没有}} \emph{méiyǒu}
\end{itemize}

\textbf{Listen.} \emph{Nǐ yǒu méiyǒu shǒujī?}

\textbf{Speak.} Answer it with \zh{有}, then with \zh{没有}, then say what your room has.

\emph{Write it:} \zh{没} — \zh{氵} first, then the small roof, then the crossing pair that opens \zh{对}
\end{grammarlens}

\subsection*{Your turn}
\begin{itemize}
\raggedright
  \item \emph{Say it:} \emph{nǐ yǒu méiyǒu yàoshi} — Have you got a key?
  \item \emph{Say it:} \zh{没有} — \emph{méiyǒu} — No.
  \item \emph{Say it:} \emph{kāi dēng} — Switch on the light.
  \item \emph{Say it:} the whole exchange, both parts
  \item \emph{Write it:} \zh{我没有手机}
\end{itemize}

\begin{tcolorbox}[breakable,colback=teal!4,colframe=teal!35!black,title={Before you move on}]
Say you don't have a phone. (\textbf{\zh{我没有手机}}.) Say you didn't go. (\textbf{\zh{我没去}}.) Say there is milk in the fridge. (\emph{bīngxiāng lǐ yǒu niúnǎi}.)
\end{tcolorbox}
